\chapter{CONCLUSION AND RECOMMENDATION}

\label{ch:conclusion}

\section{Summary of Contributions}

This thesis asked how a web-based planner can represent a university program, a student's record
and an institution's resource constraints precisely enough to be useful, without claiming more than
the available evidence supports. The artifact answers that question with seven contributions.

\begin{enumerate}
  \item \textbf{A program-aware curriculum model.} Courses are global; membership, placement and
        prerequisite edges are per program. Repeated elective placements are preserved as ordered
        placement rows, and an uncoded requirement is a first-class entity rather than a placeholder
        course. The model therefore supports several programs without duplicating shared courses
        and without unioning their prerequisite graphs.

  \item \textbf{An account scope that is enforced transactionally.} Authentication uses an
        {\tt httpOnly} cookie, the server re-reads the current role and program context on every
        request, and every mutation may carry an expected scope that is compared against the values
        read inside the writing transaction. A stale client cannot write into a scope the user no
        longer occupies, and the mechanism allowed the client to adopt it one screen at a time.

  \item \textbf{Server-authoritative progress integrity.} Completion validates mandatory
        prerequisites, and uncompletion removes direct and transitive dependents inside the same
        serializable transaction. A duplicated rule array that had disagreed with the database was
        removed, so the database is the single source of prerequisite truth.

  \item \textbf{Immutable grade evidence and a defensible average.} Grade attempts append instead of
        overwriting, carry an owner-scoped request key that makes an exact retry idempotent, and feed
        a credit-weighted average computed from the highest score per course in decimal arithmetic.
        Per-semester averages were added as a separate, narrower claim that never infers a term from
        a timestamp or a saved plan.

  \item \textbf{Difficulty estimates that report their own confidence.} Completion-gated ratings are
        combined with a program-level prior through Bayesian shrinkage, and the vote count is shown
        beside every estimate, so a course with no votes is visibly a course with no evidence.

  \item \textbf{A versioned, explicit administrative simulation.} Resource settings, cohort demand,
        capacity diagnostics and two allocation models are implemented behind frozen contracts with
        immutable captures and private per-student outcomes, executed by an explicit single-job
        command or protected action rather than an unattended daemon.

  \item \textbf{An evidence discipline.} New tests are observed failing against the code they
        replace; the server suite runs against a real database; one reproducible continuous-
        integration gate blocks every change; and the interface was audited in a real browser rather
        than assumed correct.
\end{enumerate}

The most transferable of these is the sixth. A planning tool for an unverified catalogue is a tool
whose value depends entirely on how candidly it represents its own uncertainty, and that candour has
to be built into the schema, the contracts and the copy rather than added as a disclaimer.

\section{Limitations}

\begin{enumerate}
  \item \textbf{No institutional validation.} The reference program is imported from a published
        page and is deliberately not asserted to match the signed policy of a cohort. Its credit
        total is unknown, its elective rules are unresolved, and the final-year branch is a declared
        reference rule rather than a verified regulation.
  \item \textbf{Three programs transcribed, two imported.} The model supports several programs and
        the administrative simulation exercises that shape. The programs beyond the original reference
        were analysed rather than assumed: each was retrieved from its signed document, transcribed
        and checked, and two of the three are imported.

        \textbf{Data Science is imported} as a reference cohort. Its fifty courses, fifty
        placements, thirty-two prerequisite edges and blocks of 38, 27, 32 and 16 reach the
        document's own total of 124 once its eleven elective credits are split into a seven-credit
        pool and a four-credit free elective. Two of its cells contradicted the rest of the document,
        and the transcription is committed in two forms so the contradiction stayed visible until the
        independence of cohorts resolved it: the 2025 appendix replaces the codes, and a 2025 student
        follows the 2025 plan. Importing assigns nobody, so the program is available without any
        account being moved onto it.

        \textbf{Information Technology Network Engineering is imported}, at 150 credits. Its blocks
        reconcile exactly -- 45, 35, 32 plus 12 elective, 6 and 20 -- and the construct that once
        blocked it is now built and exercised. Four of its requirements read ``IT116IU or
        IT149IU'', and \texttt{IT149IU} belongs to another program; the logic laboratory takes
        \texttt{EE054IU} as its alternative. Such an edge names no member, is excluded from the
        dependency graph, and is satisfied by a completion of that course in the student's own
        records -- a rule applied in the completion cascade, the import path, the progress summary
        and the saved-plan review, and published in the reference view rather than silently
        dropping it. Whether a completion earned under a different program should satisfy it remains
        a transfer-credit question, stated rather than decided.

        \textbf{Computer Engineering is imported} as a fourth structure, also at 150. Its blocks
        reconcile, and its schedule prints ten blocks while the document claims nine semesters --- a
        contradiction only if blocks are read as terms. They group courses by kind, so the ten blocks
        are ten separations rather than ten terms, and a student sequences their own plan subject to
        one constraint: the special study precedes the thesis and the thesis occupies the final
        semester alongside whatever remains. That rule, rather than the tables, determines the
        placements this manifest records, and it is stated as the rule's consequence rather than
        presented as a fact the document contains.

  \item \textbf{No calendar validity.} Offerings, meeting times, rooms, laboratory suitability and
        staff qualification were unavailable, so a placement in a semester is a catalogue fact and
        not evidence that the course is offered in a given term. The per-semester grade labels are
        self-recorded metadata for the same reason.
  \item \textbf{Assigned contexts are read-only.} A student whose account holds a program reference
        cannot complete courses through the map, because doing so needs a context-keyed client cache
        and a policy for transferring elective claims that were not built. The interface states the
        boundary instead of hiding it.
  \item \textbf{Simulated resources and allocations.} Every administrative number derives from
        configuration a user typed into a form. The allocation models are greedy heuristics, not
        optimisers, and their outputs are reproducible rather than correct.
  \item \textbf{Engineering evidence only.} No students were recruited and no planning outcome was
        measured. The evaluation establishes that the artifact enforces the invariants it declares,
        is reproducible, and remains operable with a keyboard and a screen reader; it does not
        establish that the recommendations are academically effective.
  \item \textbf{Not deployed.} A production start-up refuses to run without an explicit frontend
        origin, which is a configuration preflight rather than a deployment. Hosting, certificate
        and cookie topology, secret provisioning and operational recovery remain undecided.
  \item \textbf{An open stability finding.} The server suite occasionally produces an authentication
        refusal where a business-level refusal is expected, and occasionally fails a suite that
        passes in isolation. The cause is unestablished and is recorded rather than suppressed.
\end{enumerate}

\section{Future Work}

\begin{enumerate}
  \item \textbf{Reconcile the remaining programs from signed sources.} The blocking work is
        documentary, not technical, and most of it turned out to be a rule rather than a question:
        \emph{cohorts are independent}. A student admitted in a year follows that year's plan for
        their whole degree, so a later curriculum supersedes nothing, and a prior cohort's rule is
        never an alternative reading for a later one. Where the 2024 and 2025 Data Science documents
        differ, both are correct, for different students.

        That rule closes the questions rather than multiplying them. The 2025 appendix replaces
        \texttt{MA026IU} with \texttt{MA036IU} and \texttt{IT154IU} with \texttt{MA033IU}, and the
        printed \texttt{T149IU} is a slip for \texttt{IT149IU}. What remains open is narrower and
        concerns Information Technology alone: whether the code printed \texttt{IT091I} is
        \texttt{IT091IU}, whether a prerequisite given as ``Physics 1'' means \texttt{PH013IU}, and
        how its three unnamed elective slots relate to the twelve self-chosen credits its block
        allocates. A cohort must not be activated by an artefact alone.
  \item \textbf{Encode offerings and the academic calendar.} With timetable data, the reference
        year and semester could become a validated constraint rather than a label, and the
        allocation models could reason about actual meeting times and rooms.
  \item \textbf{Enable editing inside an assigned context.} This requires context-keyed browser
        caches, a claim-transfer policy for elective groups, and a decision about what happens to
        records that are not members of the newly assigned program.
  \item \textbf{Replace the greedy allocator with an optimiser} and compare the two on the same
        fixtures, so that the cost of the heuristic can be quantified rather than asserted.
  \item \textbf{Deploy and operate.} Choose a hosting topology, provision secrets securely, verify
        the deployed flows and rehearse recovery, then re-run the continuous-integration gate
        against the deployed configuration.
  \item \textbf{Evaluate with students.} A study with a control condition would be the appropriate
        way to test whether recommendations change planning decisions or outcomes. The artifact's
        existing instrumentation, which records what was recommended and what was chosen, is the
        starting point.
\end{enumerate}

\section{Recommendations}

\textbf{For the school.} Publish the curriculum in a machine-readable form with an explicit version
and an explicit effective cohort. Most of the difficulty in this work came from reconstructing
intent from a page whose credit totals, course codes and elective rules had drifted between
versions; a maintained source would remove that class of problem for every downstream tool.

\textbf{For future developers.} Two practices carried more weight than any technology choice. The
first is enforcing scope inside the transaction that writes, which turned a whole class of stale-
client bugs into conflicts. The second is treating absence as a value -- nullable credit totals,
unknown seats, unresolved requirements -- because a system that can say ``unknown'' is much harder
to mislead than one that must always produce a number.

\textbf{For students.} The artifact is most honest when used as a planning aid rather than an
authority: it can show which prerequisites are unmet, what a semester would cost in credits, and how
confident the difficulty estimates are, but it cannot tell you what the school will offer, and it
does not pretend to.
